\section{Spielleiterinformationen}

\subsection{Die Umgebung \texttt{gm}}

Die Umgebung \verb|gm| dient dazu, Informationen hervorzuheben, die ausschließlich für den Spielleiter bestimmt sind.

Sie erzeugt einen vertikalen Balken am linken Rand und hebt den Inhalt dadurch deutlich vom eigentlichen Abenteuertext ab.

\subsubsection*{Syntax}

\begin{Verbatim}
\begin{gm}
...
\end{gm}
\end{Verbatim}

\subsubsection*{Beispiel}

\paragraph{Quellcode}

\begin{Verbatim}
\begin{gm}
Die Räuber beobachten die Gruppe bereits
seit mehreren Stunden.
\end{gm}
\end{Verbatim}

\paragraph{Ergebnis}

\begin{gm}
Die Räuber beobachten die Gruppe bereits
seit mehreren Stunden.
\end{gm}

\newpage

\subsection{Der Befehl \texttt{\textbackslash gmnote}}

Dieser Befehl ist eine Kurzform der Umgebung \verb|gm|.

Er eignet sich besonders für kurze Hinweise.

\subsubsection*{Syntax}

\begin{Verbatim}
\gmnote{Text}
\end{Verbatim}

\subsubsection*{Beispiel}

\paragraph{Quellcode}

\begin{Verbatim}
\gmnote{Der Wirt lügt über seine Vergangenheit.}
\end{Verbatim}

\paragraph{Ergebnis}

\gmnote{Der Wirt lügt über seine Vergangenheit.}
